\newcommand{\reportAuthor} {%
  Mouad \textsc{Ettahiri}%
}

\newcommand{\reportTitle} {%
  Conception et développement d'une application pour l'automatisation de la gestion des concours de recrutement%
}

\newcommand{\reportSubject} {%
  Conception et développement d'une application pour l'automatisation de la gestion des concours de recrutement%
}

\newcommand{\dateSoutenance} {%
  Septembre 2026%
}

\newcommand{\studyDepartment} {%
  Ingénierie Informatique et Réseaux option MIAGE%
}

\newcommand{\ENETCOM} {%
  École Marocaine des Sciences de l'Ingénieur -- Rabat%
}

\newcommand{\codePFE} {%
}

\newcommand{\juryPresident} {%
  Pr. Imane \textsc{Hilal}%
}
\newcommand{\juryPresidentDesc} {%
  Tuteur de l'école%
}

\newcommand{\juryMemberOne} {%
  M. Tarik \textsc{Lakhbizi}%
}
\newcommand{\juryMemberOneDesc} {%
  Tuteur de stage%
}

\newcommand{\juryMemberTwo} {%
}
\newcommand{\juryMemberTwoDesc} {%
}

\newcommand{\specialcell}[1]{%
  \begin{tabularx}{\textwidth}{@{}X@{}}#1\end{tabularx}%
}

%%%%%%%%%%%%%%%%%%%%%%%%%%%%%%%%%%%%%%%%%%%%%%%%%%%%%%%
% Add your own commands here
%%%%%%%%%%%%%%%%%%%%%%%%%%%%%%%%%%%%%%%%%%%%%%%%%%%%%%%
\newcommand{\MyCommand} {%
  Does nothing really%
}

% Diagramme principal : page dédiée, portrait (pas de rotation)
\newcommand{\fullpagefig}[3]{%
  \ifdim\pagetotal>0pt\clearpage\fi
  \begin{figure}[H]
    \centering
    \includegraphics[width=\textwidth,height=0.88\textheight,keepaspectratio]{#1}%
    \caption{#2}
    \label{#3}
  \end{figure}
  \clearpage
}

% Sous-diagramme (a, b, c) : compact, plusieurs par page si possible
\newcommand{\subseqfig}[3]{%
  \begin{figure}[H]
    \centering
    \includegraphics[width=\textwidth,height=0.34\textheight,keepaspectratio]{#1}%
    \caption{#2}
    \label{#3}
  \end{figure}%
}

% Conservé si un diagramme large devait revenir en paysage
\newcommand{\landscapefig}[3]{%
  \begin{landscape}
    \begin{figure}[H]
      \centering
      \includegraphics[width=0.97\linewidth,height=0.78\textheight,keepaspectratio]{#1}%
      \caption{#2}
      \label{#3}
    \end{figure}
  \end{landscape}
}
